\clearpage
\section*{Anexo 4.1-U --- Prueba de entrega, guía y acuses}
\phantomsection
\label{anx:U}
\addcontentsline{toc}{section}{Anexo 4.1-U --- Prueba de entrega, guía y acuses}
La solución propone tres registros separados: documento tributario emitido por ERP, evidencia operacional capturada por M6 y acuse recibido por el canal del destinatario. El SII exige portar la representación gráfica o impresa del DTE durante traslado; su acuse técnico no acredita por sí solo recepción de mercaderías (Servicio de Impuestos Internos [SII], s.~f.-a). La evidencia móvil tampoco se transforma automáticamente en el recibo legal del destinatario.

\subsection*{Mecanismo de recepción propuesto}
M5 guarda folio, tipo, emisor, destinatario, versión de carga, documento y hash antes de liberar; el conductor dispone de representación gráfica o impresa conforme al procedimiento tributario. En el local, M6 registra nombre del receptor, identificación según política validada, lugar, hora, cantidades aceptadas/rechazadas y firma manuscrita capturada o alternativa operacional QR/fotografía. La cola cifrada conserva el objeto y su hash; la sincronización vincula la evidencia a entrega y guía.

Para un receptor electrónico habilitado se propone recibir su acuse por el mecanismo tributario del ERP o el canal EDI acordado, conservar el XML y su validación e identificar al firmante o emisor autorizado. Para el receptor al que corresponde constancia en representación impresa, se obtiene esa constancia y se conserva copia digital vinculada; no se obliga a una persona natural a poseer firma electrónica avanzada. La diferencia se fundamenta en la interpretación publicada por el SII sobre receptores electrónicos y no obligados (SII, 2018) y en su formato de recibo electrónico (SII, s.~f.-b). Tributación valida el perfil aplicable de cada cliente antes de habilitarlo.

\begin{tablalafrox}{Estados de evidencia y consecuencias}{tab:estados-documentales}{F{0.23} F{0.29} Y}{\cab{Hecho} & \cab{Evidencia conservada} & \cab{Consecuencia lógica}}
Guía disponible & Folio, documento y hash del ERP. & M5 comprueba correspondencia con carga vigente. \\
Salida liberada & Usuario, carga, guía y hora. & M6 recibe el viaje; no modifica el documento. \\
Recepción completa & Cantidades y POD íntegro. & M7 concilia; acuse tributario sigue su propio estado. \\
Recepción parcial & Líneas rechazadas, lotes y causal. & M8 registra devolución; ERP decide ajuste documental. \\
Negativa a firmar & Nombre si disponible, causal y constancia alternativa. & Entrega en excepción; no marcar acuse legal obtenido. \\
Acuse tributario recibido & Origen, XML/constancia, fecha y validación. & Vincular al DTE sin reemplazar el POD. \\
Resultado incierto & Clave original y último estado conocido. & Consultar antes de reenviar o solicitar otro documento. \\
\end{tablalafrox}
{\footnotesize Fuente: diseño lógico M5--M8, RT-16.14 y publicaciones del SII citadas.\par}

AL-POD-01 ensaya recepción completa, parcial, persona sustituta, negativa a firmar, falta de señal y objeto alterado. Se acepta cuando cada hecho conserva autor y origen, la guía precede a la salida, el documento no cambia por reenvío y una evidencia operacional insuficiente no se declara acuse tributario. La FAQ del SII sobre contingencia es orientativa y no equivale a autorización para este contribuyente (SII, s.~f.-c); el procedimiento de emisión bajo falla se cierra por AL-DTE-01.
